%%%%%%%%%%%%%%%%%%%%%%%%%%%%%%%%%%%%%%%%%%%%%%%%%%%%%%%%%%%%%%%%%%%%%%
% Lớp 9 | Chương IV. Hệ thức lượng trong tam giác vuông | Tỉ số lượng giác của góc nhọn
% Dạng 9H4?11-2: Tỉ số lượng giác của hai góc phụ nhau
%%%%%%%%%%%%%%%%%%%%%%%%%%%%%%%%%%%%%%%%%%%%%%%%%%%%%%%%%%%%%%%%%%%%%%


% nguồn: dan_tu_AI (nhập 17/11/2026)
\begin{ex}%[9H4N11-2]%[ma:7bb03f24]
	Khẳng định nào sau đây đúng?
	\choice
	{$\sin 40^\circ=\sin 50^\circ$}
	{$\sin 40^\circ=\cos 40^\circ$}
	{$\tan 40^\circ=\tan 50^\circ$}
	{\True $\sin 40^\circ=\cos 50^\circ$}
	\loigiai{
		Hai góc $40^\circ$ và $50^\circ$ phụ nhau ($40^\circ+50^\circ=90^\circ$) nên $\sin 40^\circ=\cos 50^\circ$.
	}
\end{ex}

% nguồn: dan_tu_AI (nhập 17/11/2026)
\begin{ex}%[9H4H11-2]%[ma:c3945dee]
	Các khẳng định sau đúng hay sai?
	\choiceTF
	{\True $\sin 25^\circ=\cos 65^\circ$}
	{$\tan 50^\circ=\cot 50^\circ$}
	{\True $\tan 35^\circ\cdot\tan 55^\circ=1$}
	{$\cos 15^\circ-\sin 75^\circ+\tan 20^\circ\cdot\tan 70^\circ=2$}
	\loigiai{
		\begin{itemchoice}
			\itemch Đúng. $25^\circ+65^\circ=90^\circ$ nên $\sin 25^\circ=\cos 65^\circ$.
			\itemch Sai. $\tan 50^\circ=\cot 40^\circ$; còn $\cot 50^\circ=\tan 40^\circ<\tan 45^\circ=1<\tan 50^\circ$.
			\itemch Đúng. $\tan 55^\circ=\cot 35^\circ$ nên $\tan 35^\circ\cdot\tan 55^\circ=\tan 35^\circ\cdot\cot 35^\circ=1$.
			\itemch Sai. $\cos 15^\circ=\sin 75^\circ$ và $\tan 70^\circ=\cot 20^\circ$ nên biểu thức bằng $0+1=1$.
		\end{itemchoice}
	}
\end{ex}

% nguồn: dan_tu_AI (nhập 17/11/2026)
\begin{bt}%[9H4H11-2]%[ma:155eb949]
	Tính giá trị của biểu thức $A=\dfrac{\sin 32^\circ}{\cos 58^\circ}+\tan 17^\circ\cdot\tan 73^\circ-\cos^2 45^\circ$.
	\shortans{1,5}
	\loigiai{
		$32^\circ+58^\circ=90^\circ$ nên $\sin 32^\circ=\cos 58^\circ$, suy ra $\dfrac{\sin 32^\circ}{\cos 58^\circ}=1$.\\
		$17^\circ+73^\circ=90^\circ$ nên $\tan 73^\circ=\cot 17^\circ$, suy ra $\tan 17^\circ\cdot\tan 73^\circ=1$.\\
		$\cos^2 45^\circ=\left(\dfrac{\sqrt2}{2}\right)^2=\dfrac12$.\\
		Vậy $A=1+1-\dfrac12=1{,}5$.
	}
\end{bt}

% nguồn: dan_tu_AI (nhập 17/11/2026)
\begin{bt}%[9H4C11-2]%[ma:9b14ef72]
	Tính giá trị của tổng
	\[ S=\sin^2 1^\circ+\sin^2 2^\circ+\sin^2 3^\circ+\cdots+\sin^2 88^\circ+\sin^2 89^\circ. \]
	\shortans{44,5}
	\loigiai{
		Với $1\le k\le 44$, hai góc $k^\circ$ và $(90-k)^\circ$ phụ nhau nên $\sin(90-k)^\circ=\cos k^\circ$, do đó
		\[ \sin^2 k^\circ+\sin^2(90-k)^\circ=\sin^2 k^\circ+\cos^2 k^\circ=1. \]
		Tổng $S$ gồm $44$ cặp như vậy và số hạng $\sin^2 45^\circ=\dfrac12$.\\
		Vậy $S=44\cdot 1+\dfrac12=44{,}5$.
	}
\end{bt}
